% MAIN SNAPSHOT — LSPV/RMV::SYNTHESIS-0026::2026-10-12T20:46:09Z
% !TEX program = xelatex
\documentclass[11pt]{article}
\usepackage[a4paper,margin=0.88in]{geometry}
\usepackage{fontspec}
\usepackage{amsmath,amssymb,amsthm,booktabs,tabularx,array,xurl,hyperref,microtype}
\hypersetup{colorlinks=true,linkcolor=blue,urlcolor=blue,citecolor=blue}
\setlength{\emergencystretch}{3em}
\newtheorem{theorem}{Theorem}
\newtheorem{corollary}{Corollary}
\newcommand{\EXACT}{\textbf{PROVED/EXACT}}
\newcommand{\KNOWN}{\textbf{KNOWN THEOREM}}
\newcommand{\COMP}{\textbf{COMPUTATIONAL}}
\newcommand{\COND}{\textbf{CONDITIONAL}}
\newcommand{\CAND}{\textbf{CANDIDATE}}
\newcommand{\OPEN}{\textbf{OPEN}}
\newcommand{\RET}{\textbf{FALSIFIED/RETIRED}}
\newcommand{\same}{\textnormal{\textsc{same\_realization}}}
\title{Universal-Core Synthesis: Finite Fault Domains Do Not Cover All Polynomial-Time Deciders}
\author{LSPV-5 / MAIN synthesis}
\date{13 October 2026}

\begin{document}
\sloppy
\maketitle

\section*{Executive status}
\textbf{P versus NP remains open.}  The run proves a negative universality
theorem for the most concrete repair suggested by the preceding black-box
testing result.  Complete finite-state conformance suites are valid relative
to an explicitly bounded fault domain.  They do not yield a universal bridge
for all uniform polynomial-time SAT deciders: the class has no common
polynomial exponent, and finite exceptional behaviour defeats every exact
per-layer finite intensional cap.  Quotienting implementations extensionally
repairs finiteness only by forgetting the internal trace/resource property the
test was supposed to certify.

\paragraph{Strict status.}
\textsf{FAULT-DOMAIN-UNIVERSALITY-TRILEMMA1} is \EXACT{} as an elementary
ordinary-machine theorem.  Classical FSM conformance completeness inside a
declared state-bounded fault domain is a \KNOWN{} result.  The bounded hostile
audit is \COMP{}.  The proposed bridge
\textsf{FINITE-FAULT-DOMAIN-COMPLETENESS-AS-UNIVERSAL-BRIDGE1} is \RET{}.
Program-relative or compiler-image conformance remains a \CAND{}, but it does
not quantify over arbitrary SAT deciders.  The all-decider clean-core lower
bound remains \OPEN{}.

\section*{Recovery and canonical preflight}
The canonical repository contains the verified predecessor
\path{research/agents/lspv-4/2026-10-13_0456_LSPV-4_PHASEBRIDGE1.tex}, with
the required header and sentinel.  It also contains the later verified MAIN
snapshot
\path{research/main/2026-10-12_1547_MAIN_SYNTHESIS-0025.tex}.  Thus the older
recovery premise naming SYNTHESIS-0023 is stale state, not a current missing
write.  No unavailable report body was reconstructed and no duplicate recovery
file was created.

\section*{Exact clean core and target}
Retain the root manuscript's fixed NP-complete clean core
\(\mathcal H_{2,2}\): duplicate-free proper Monotone
\(3\mathrm{SAT}\text{-}(2,2)\) under the canonical indexed encoding.  If
\(F\in\mathcal H_{2,2}\), write
\[
 N_F=|\operatorname{enc}(F)|=\Theta(n\log n).
\]
For a partial assignment \(\alpha\), semantic compatibility is
\[
 H(F,\alpha)=\{\beta:\alpha\cup\beta\models F\},\qquad
 h(F,\alpha)={\bf 1}[H(F,\alpha)\ne\varnothing].
\]
The exact unrestricted target is
\[
 \forall A\in\operatorname{Dec}_{\rm uniform}(\mathcal H_{2,2})\;
 \forall d\;\forall N_0\;\exists F\in\mathcal H_{2,2}:
 N_F\ge N_0\ \land\ T_A(F)>N_F^d.                 \tag{UCORE-TIME1}
\]
Since the core is NP-complete, this target is separation-strength.  No report
proves it.  The point of the present theorem is narrower: it rules out a finite
conformance fault domain as a weaker route around that quantifier.

\section*{Obstacle and cross-domain transfer checked first}
The two current obstacles are:
\begin{enumerate}
\item black-box testing cannot certify a trace property at an unqueried input;
\item a complete conformance theorem needs a fault domain that actually
contains every implementation quantified over by the desired theorem.
\end{enumerate}
Chow's FSM testing construction and its descendants solve the second problem
only after fixing a finite specification and a bounded implementation class.
Modern completeness analyses commonly parameterize the domain by an upper
bound on additional implementation states; the resulting suites grow
exponentially with that allowance.  This is a \emph{directly transferable
theorem} for a declared finite transition-system model.  Treating its state
bound as a theorem about arbitrary ordinary polynomial-time programs is a
\emph{failed transfer}: the objects and quantifiers differ.  VLSI controller
state bounds and finite verification domains are useful structural analogies,
but they likewise require an explicit implementation contract.

\section*{Theorem: FAULT-DOMAIN-UNIVERSALITY-TRILEMMA1}
Let \(L\subseteq\{0,1\}^*\), and suppose \(A\) is a uniform deterministic
polynomial-time decider for \(L\).  Let an intensional fault domain attempt to
contain all uniform polynomial-time deciders extensionally equivalent to
\(A\).

\begin{theorem}
The following three statements hold.
\begin{enumerate}
\item For every fixed \(c\ge0\), there is an extensionally equivalent uniform
polynomial-time decider \(D_c(A)\) whose realized execution on every input of
length \(N\ge2\) visits at least \(N^{c+1}\) distinct counter configurations
before running \(A\).  Hence a domain capped by \(N^c\) configurations or
steps does not contain all polynomial-time equivalent deciders.
\item For every fixed length \(N_0\) and every finite cap \(B(N_0)\), there is
an extensionally equivalent uniform polynomial-time decider
\(D_{N_0,B}(A)\) whose execution on every length-\(N_0\) input visits more
than \(B(N_0)\) counter configurations before running \(A\).  Hence no exact
finite cap assigned separately to each layer contains all such deciders.
\item If the fault domain instead identifies all implementations having the
same one-bit behaviour on each input, then all correct deciders of \(L\) lie
in one extensional class.  Any nonconstant predicate of their internal trace,
runtime, representation, or memory layout is not well-defined on that
quotient.
\end{enumerate}
\end{theorem}

\paragraph{Proof.}
For (1), on input \(x\) of length \(N\), initialize a counter to zero, execute
exactly \(N^{c+1}\) increments, erase or ignore the counter, then run \(A(x)\)
and return its bit.  The configurations with different counter values are
distinct.  The wrapper is uniform and runs in polynomial time, while
\(N^{c+1}>N^c\) for \(N\ge2\).

For (2), hard-code the finite integers \(N_0\) and
\(C=B(N_0)+1\).  On inputs of length \(N_0\), perform \(C\) counter
increments; on all other lengths skip them; then run \(A\).  The exceptional
cost is a finite machine constant, so the resulting machine is still uniform
polynomial time.  At the selected layer it exceeds the cap.

For (3), extensional equality identifies programs precisely by the function
they compute.  Every correct decider computes the characteristic function
\(\chi_L\), yet COPY, DELAY, recomputation, interpretation, and syntax-padding
wrappers can change the named intensional property without changing
\(\chi_L\).  Therefore that property cannot descend to the quotient.
\hfill\(\square\)

\paragraph{Quantifier guard.}
There is no claim that one fixed wrapper defeats every possible cap at every
length.  Part (1) responds to each fixed exponent \(c\); part (2) responds to
each pair \((N_0,B(N_0))\).  This is exactly enough to refute one universal
fixed-polynomial cap and every exact finite layer cap.  It does not refute a
correctly declared restricted fault domain.

\section*{Consequence for the proposed bridge}
The preceding BBTEST1 result showed that a finite query audit cannot certify a
universal trace property outside its queried layer points; its smallest repair
was a complete fault domain or a density theorem.  The theorem above closes
the first repair at arbitrary-program scope:
\begin{itemize}
\item a bounded intensional domain excludes correct polynomial-time wrappers;
\item an unbounded domain no longer supplies the finite completeness premise;
\item an extensional quotient includes the wrappers only by erasing the
resource property;
\item a compiler-defined domain can be complete only for the compiler image,
not for arbitrary deciders without another universality theorem.
\end{itemize}
The density alternative remains \OPEN{}, but it must be a theorem about the
same realized executions, not a sampling heuristic.

\section*{Bounded reproducible audit}
The script
\path{experiments/review/fault_domain_universality_trilemma_audit.py}
performs bounded checks for \(N=2,\ldots,8\), cap exponents
\(c=0,\ldots,4\), and exact layer caps
\(0,1,2,3,7,15,31,63\).  It tests six extensionally equivalent
presentations (DIRECT, COPY, DEAD, DELAY, INTERP, RECOMP) over 2SAT, Horn,
affine/Tseitin, bounded-treewidth, independent-product, padded, redundant, and
alternate-encoding controls.  Each row explicitly enumerates its distinct
counter values once; hostile executions use exact counter arithmetic rather
than replaying the same delay.

The run completed in \(0.30\) wall-clock seconds, with \(368{,}318\)
assertions, \(121{,}920\) hostile executions, 35 exponent rows, 56 exact-layer
rows, 600 distinct syntactic presentations, 20 extensional truth-table
classes, and zero failures.  The result is
\path{results/raw/review/2026-10-13_fault_domain_universality_trilemma.json}.
Script SHA-256 is
\texttt{10957e4ebf4c29055787e8dd94c1505ba621c1a8658c13e0cac90edbbfceb582};
result SHA-256 is
\texttt{cfeca317edcec654a90a3f7a5a743bd9d23a498d448399e1467029fe0a4ba3c0}.
The finite audit is not an asymptotic proof; the theorem above supplies the
symbolic argument.

\section*{Latest-peer reconciliation}
\begin{center}\small
\begin{tabularx}{\textwidth}{p{0.18\textwidth}Xp{0.18\textwidth}}
\toprule
Artifact & Verified contribution and current consequence & Status here\\
\midrule
LSPV-1 QNORM1 & The fully quantified fixed-core time lower bound is exactly
separation-strength; there is no weaker conclusion hidden by notation. &
Adopted\\
LSPV-2 BBTEST1 & A finite black-box audit misses a one-input trace mutation
unless a complete fault domain or density theorem is supplied. & Strengthened:
finite-domain repair now retired at arbitrary-program scope\\
LSPV-3 LOGCOF1 & A logarithmic cofactor table lets a wrapper perform the
semantic work before the claimed late decision phase. & Adopted hostile
control\\
LSPV-4 PHASEBRIDGE1 & With \(k=O(\log N)\) withheld bits, a solver may
precompute all cofactors before source coverage while preserving polynomial
time. & Adopted; phase-local charge retired\\
\bottomrule
\end{tabularx}
\end{center}
Together these reports force the target back to whole-execution worst-case
cost with arbitrary polynomial exponent and exact \same{} accounting.

\section*{Faithful comparison with actual main4 and main5}
The current \path{main4.tex} is the ARC achievable-region manuscript and
\path{main5.tex} is the WRAPPER construction manuscript.  They contain useful
construction-specific controllers, reductions, and charged accounting.  They
do not prove that every uniform polynomial-time SAT decider belongs to one
finite implementation fault domain, exposes \(H(F,\alpha)\), or realizes a
particular ARC/WRAPPER trace.  The present theorem does not alter their valid
restricted statements and does not silently reinterpret either manuscript as
an all-algorithm lower bound.

\section*{Barrier, model, and SAME\_REALIZATION audit}
\paragraph{Relativization and algebrization.}
The delay wrappers work unchanged relative to an oracle and with algebraic
oracle access.  The theorem is therefore a negative guard, not a
nonrelativizing or nonalgebrizing hinge.

\paragraph{Natural proofs and nonuniformity.}
No large constructive circuit property is produced.  Circuit, branching
program, proof-system, VLSI, and FSM lower bounds stay in their declared
models.  The theorem concerns uniform machines individually, but their
polynomial exponents and finite constants are not uniformly bounded over all
machines.

\paragraph{Same realization and all costs.}
Every wrapper execution is a separate realized program.  Its length check,
counter initialization and increments, invocation of \(A\), cleanup, and
one-bit output are charged to that wrapper.  No semantic state count is used to
infer runtime: the construction actually traverses the counter
configurations.  Conversely, a conformance tester's generation, execution,
comparison, extraction, and reconstruction costs belong to the tester, not to
the tested SAT run.

\paragraph{One-bit output.}
The wrapper returns exactly the bit returned by \(A\).  It need not expose a
witness, compatibility table, proof, or trace certificate.  Extensional
correctness therefore cannot select one intensional representative.

\section*{Post-quantum technique ledger}
Lossiness is an \emph{exact object transfer} only for the extensional quotient:
many program presentations collide on the same truth table.  That lossiness is
algorithmically exploitable as a shortcut and erases provenance; it is not a
hardness source.  Schur-square, quadratic-relation, Hilbert/syzygy, and
cryptographic binding techniques are \emph{failed transfers} for this barrier
because no code or unconditional binding object is present.  Worst-case to
average-case extraction remains an \emph{analogy/candidate} requiring a new
density theorem and an audit of known NP worst-case/average-case barriers.
Hash and random-oracle methods remain inadmissible without removing their
assumptions.

\section*{P=NP counterfactual audit}
Assume the fixed clean core has a uniform polynomial-time decider \(A\).  The
wrappers above are then also polynomial-time deciders of the same core.  This
produces no contradiction; it merely demonstrates why a universal cap cannot
be inferred from the word ``polynomial.''  To contradict the assumption one
still needs either an unconditional lower bound for the same general machine
model or a transfer placing every such decider into one restricted model with
a model-matched lower bound.  Assuming the required cap or restricted-model
placement would smuggle the missing separation into the premise.

\section*{Proof spine and highest-value target}
\[
\begin{array}{c}
\text{arbitrary uniform polynomial clean-core decider }A\\
\Downarrow\quad\text{\OPEN: invariant same-realization transfer}\\
\text{whole-execution charged object in one model}\\
\Downarrow\quad\text{\OPEN: unconditional lower bound in that model}\\
\text{contradiction.}
\end{array}
\]
The highest-value target is now
\textsf{DENSITY-OR-DIRECT-TOTAL-COST1}: either prove an all-decider density
theorem for a trace/resource violation that survives COPY, DELAY, RECOMP,
INTERP, affine elimination, bounded treewidth, padding, and source-first
cofactor precomputation, or prove UCORE-TIME1 directly.  A finite test suite
may be used only after its fault domain is declared and proved to contain the
same implementations under study.

\section*{Corpus reviewed}
The run surveyed the canonical tree; read root \path{main.tex},
\path{AGENTS.md}, \path{CODEX.md}, the repository research skill, all task
specifications, actual \path{main4.tex} and \path{main5.tex}, SYNTHESIS-0025,
the newest verified LSPV-1--4 reports, and materially relevant older wrapper,
trace, compiler, and quotient reports.  It also reviewed the HEC, VSC-HEC,
G-PPC-HEC, CRS-PPC, and CCSPC explorations, governance notes, and the BBTEST1
experiment/result artifacts.  Repository search found no earlier canonical
proof of this exact fault-domain trilemma.

\section*{Duplication / contradiction audit}
The result does not duplicate BBTEST1: that report proves a black-box query
indistinguishability bound and proposes a complete fault domain as one repair;
this report tests and retires that repair only at arbitrary-program scope.
It agrees with QNORM1's arbitrary-exponent guard and with earlier wrapper-orbit
results.  There is no contradiction with finite-state conformance theory,
because those theorems assume the bounded fault domain denied here for the
universal SAT quantifier.

\section*{What changed relative to root main}
One open repair is now classified more sharply.  ``Use a complete finite fault
domain'' is no longer an admissible universal bridge unless a separate theorem
places every arbitrary polynomial-time SAT decider in that domain.  The open
team target is correspondingly narrowed to a density theorem over same
realizations or the direct whole-execution clean-core lower bound.

\section*{Adopted / rejected / retired prior components}
\begin{itemize}
\item \textbf{Adopted:} BBTEST1's query indistinguishability theorem;
QNORM1's exact quantifier normal form; LOGCOF1 and PHASEBRIDGE1 as source-first
and phase-shift hostile controls; construction-specific finite conformance
testing.
\item \textbf{Rejected as transfer:} importing FSM/VLSI state bounds into
ordinary polynomial-time programs without an object correspondence.
\item \textbf{Retired:} finite-fault-domain completeness as a universal
all-decider bridge; phase-local paid compatibility; semantic multiplicity as
runtime or memory by itself.
\item \textbf{Retained open:} an all-decider density theorem,
\same{} transfer, model-matched endpoint lower bound, and UCORE-TIME1.
\end{itemize}

\section*{Stale or overstated claims found}
Any statement that ``polynomial-time implementations have polynomially many
states'' must quantify the exponent, constants, input layer, and machine
model.  There is no one fixed polynomial cap over all machines in \(P\), and
finite exceptional behaviour defeats exact per-layer caps.  Likewise,
conformance completeness without its fault-domain hypothesis, or an
extensional quotient claimed to preserve runtime provenance, is overstated.
All such universal readings are \RET{}.

\section*{Current ledger}
\begin{center}\small
\begin{tabular}{p{0.58\textwidth}p{0.22\textwidth}}
\toprule
Item & Status\\
\midrule
Fault-domain universality trilemma & \EXACT{}\\
FSM conformance within declared bounded domain & \KNOWN{}\\
Finite hostile audit & \COMP{}\\
Finite fault domain as all-decider bridge & \RET{}\\
Program-relative/compiler-image conformance & \CAND{}; restricted scope\\
All-decider density theorem & \OPEN{}\\
Whole-execution clean-core lower bound & \OPEN{}\\
P versus NP & \OPEN{}\\
\bottomrule
\end{tabular}
\end{center}

\section*{Peer handoff}
LSPV-1 should formulate any proposed resource with its exact machine-level
invariance and arbitrary-exponent quantifiers.  LSPV-2 should test whether a
density repair can survive a one-input mutation and delay/source-first
wrappers without assuming a sampling distribution.  LSPV-3 should attack the
density quantifier with sparse exceptional lengths and syntax/state recodings.
LSPV-4 should accept finite conformance only for a declared compiler image and
keep tester cost separate from the tested execution.  LSPV-5 should reject any
future finite test claim that omits the fault domain or silently bounds all
polynomial exponents.

\section*{Primary references checked}
\begin{itemize}
\item T. S. Chow, ``Testing software design modeled by finite-state
machines,'' \emph{IEEE Transactions on Software Engineering} SE-4(3), 1978,
178--187, \url{https://doi.org/10.1109/TSE.1978.231496}.
\item F. Vaandrager, R. Fiter\u{a}u-Bro\c{s}tean, and M. Melse,
``Completeness of FSM test suites reconsidered,'' 2024,
\url{https://arxiv.org/abs/2401.12703}.
\item F. Vaandrager et al., ``New fault domains for conformance testing of
finite state machines,'' \emph{CONCUR 2025},
\url{https://doi.org/10.4230/LIPIcs.CONCUR.2025.1}.
\item T. Baker, J. Gill, and R. Solovay, ``Relativizations of the
\(P=?NP\) question,'' \emph{SIAM Journal on Computing} 4(4), 1975,
431--442, \url{https://doi.org/10.1137/0204037}.
\item S. Aaronson and A. Wigderson, ``Algebrization: a new barrier in
complexity theory,'' \emph{ACM TOCT} 1(1), 2009,
\url{https://doi.org/10.1145/1490270.1490272}.
\item A. A. Razborov and S. Rudich, ``Natural proofs,'' \emph{JCSS} 55(1),
1997, 24--35, \url{https://doi.org/10.1006/jcss.1997.1494}.
\end{itemize}

% END_OF_VERIFIED_MAIN::SYNTHESIS-0026
\end{document}
